\textbf{Phần 1. Trắc nghiệm lựa chọn}
\setcounterpageref{ex}{0}

\begin{ex} %[2D1N03]
Trong phép tính $32 + 15 = 47$, số 47 được gọi là
\choice
{Số hạng}
{\True Tổng}
{Số bị trừ}
{Hiệu}
\loigiai{
Kết quả của phép cộng gọi là Tổng.\\
\textbf{Đáp án của câu hỏi là B}
}
\end{ex}

\begin{ex} %[2D1N03]
Trong phép tính $58 - 23 = 35$, số 58 được gọi là
\choice
{Số trừ}
{Hiệu}
{\True Số bị trừ}
{Tổng}
\loigiai{
Số đứng trước dấu trừ gọi là Số bị trừ.\\
\textbf{Đáp án của câu hỏi là C}
}
\end{ex}

\begin{ex} %[2D1H03]
Trong phép tính $58 - 23 = 35$, số 23 được gọi là
\choice
{\True Số trừ}
{Số bị trừ}
{Hiệu}
{Số hạng}
\loigiai{
Số đứng sau dấu trừ gọi là Số trừ.\\
\textbf{Đáp án của câu hỏi là A}
}
\end{ex}

\begin{ex} %[2D1H03]
Trong phép tính $58 - 23 = 35$, số 35 được gọi là
\choice
{Tổng}
{Số trừ}
{Số bị trừ}
{\True Hiệu}
\loigiai{
Kết quả phép trừ gọi là Hiệu.\\
\textbf{Đáp án của câu hỏi là D}
}
\end{ex}

\begin{ex} %[2D1H03]
Cho hai số hạng lần lượt là $24$ và $15$. Tổng của hai số đó là
\choice
{$38$}
{\True $39$}
{$40$}
{$9$}
\loigiai{
$24 + 15 = 39$.\\
\textbf{Đáp án của câu hỏi là B}
}
\end{ex}

\begin{ex} %[2D1H03]
Số bị trừ là $68$, số trừ là $24$. Hiệu là
\choice
{$42$}
{$46$}
{$92$}
{\True $44$}
\loigiai{
$68 - 24 = 44$.\\
\textbf{Đáp án của câu hỏi là D}
}
\end{ex}

\begin{ex} %[2D1V03]
Tìm tổng của số bé nhất có hai chữ số và số lớn nhất có một chữ số.
\choice
{$18$}
{$20$}
{\True $19$}
{$10$}
\loigiai{
Số bé nhất có 2 chữ số: 10. Số lớn nhất có 1 chữ số: 9. Tổng: $10 + 9 = 19$.\\
\textbf{Đáp án của câu hỏi là C}
}
\end{ex}

\begin{ex} %[2D1V03]
Tìm hiệu của số tròn chục lớn nhất có hai chữ số và $40$.
\choice
{\True $50$}
{$40$}
{$60$}
{$70$}
\loigiai{
Số tròn chục lớn nhất: 90. Hiệu: $90 - 40 = 50$.\\
\textbf{Đáp án của câu hỏi là A}
}
\end{ex}

\begin{ex} %[2D1V03]
Trong một phép trừ có số bị trừ bằng $45$ và số trừ bằng $0$. Hiệu là
\choice
{$0$}
{$54$}
{$1$}
{\True $45$}
\loigiai{
$45 - 0 = 45$.\\
\textbf{Đáp án của câu hỏi là D}
}
\end{ex}

\begin{ex} %[2D1V03]
Trong một phép trừ có số bị trừ bằng số trừ. Hiệu là
\choice
{\True $0$}
{$1$}
{Bằng số bị trừ}
{Không tính được}
\loigiai{
Hai số bằng nhau trừ cho nhau có hiệu bằng 0.\\
\textbf{Đáp án của câu hỏi là A}
}
\end{ex}

\begin{ex} %[2D1V03]
Nếu một số hạng bằng 0 thì tổng sẽ
\choice
{Bằng 0}
{\True Bằng số hạng còn lại}
{Bằng 1}
{Tăng gấp đôi}
\loigiai{
$a + 0 = a$ (bằng số hạng còn lại).\\
\textbf{Đáp án của câu hỏi là B}
}
\end{ex}

\begin{ex} %[2D1C03]
Trong một phép trừ, nếu tăng số bị trừ thêm 5 đơn vị và giữ nguyên số trừ thì hiệu sẽ
\choice
{Giảm đi 5 đơn vị}
{Không thay đổi}
{\True Tăng thêm 5 đơn vị}
{Tăng gấp đôi}
\loigiai{
Tăng số bị trừ bao nhiêu đơn vị thì hiệu tăng bấy nhiêu đơn vị.\\
\textbf{Đáp án của câu hỏi là C}
}
\end{ex}

\textbf{Phần 2. Trắc nghiệm đúng sai}
\setcounterpageref{ex}{0}

\begin{ex} %[2D1H03]
Xét tính đúng sai về tên gọi thành phần phép tính:
\choiceTF
{\True Trong phép cộng $20 + 30 = 50$, $20$ và $30$ là các số hạng.}
{\True Trong phép trừ $50 - 20 = 30$, $50$ là số bị trừ.}
{Trong phép trừ $50 - 20 = 30$, $30$ là số trừ.}
{\True Trong phép trừ $50 - 20 = 30$, $20$ là số trừ và $30$ là hiệu.}
\loigiai{
\begin{itemize}
    \item $20$ và $30$ là số hạng: Đúng.
    \item $50$ là số bị trừ: Đúng.
    \item $30$ là hiệu (không phải số trừ). Ý 3 Sai, Ý 4 Đúng.
\end{itemize}
\begin{itemize}
    \item a) \textbf{Đúng}.
    \item b) \textbf{Đúng}.
    \item c) \textbf{Sai}.
    \item d) \textbf{Đúng}.
\end{itemize}
}
\end{ex}

\begin{ex} %[2D1V03]
Cho phép tính $45 + 12 = 57$. Xét tính đúng sai:
\choiceTF
{\True Tổng là $57$.}
{\True Các số hạng là $45$ và $12$.}
{Nếu đổi chỗ hai số hạng thì tổng thay đổi.}
{\True Khi đổi chỗ các số hạng thì tổng không thay đổi ($12 + 45 = 57$).}
\loigiai{
\begin{itemize}
    \item Tổng là 57: Đúng.
    \item Số hạng 45 và 12: Đúng.
    \item Phép cộng có tính chất giao hoán nên không đổi. Ý 3 Sai, Ý 4 Đúng.
\end{itemize}
\begin{itemize}
    \item a) \textbf{Đúng}.
    \item b) \textbf{Đúng}.
    \item c) \textbf{Sai}.
    \item d) \textbf{Đúng}.
\end{itemize}
}
\end{ex}

\textbf{Phần 3. Trả lời ngắn}
\setcounterpageref{ex}{0}

\begin{ex} %[2D1H03]
Tính tổng của hai số hạng $35$ và $24$.
\shortans{59}
\loigiai{
\textbf{Đáp số:} $59$\\
$35 + 24 = 59$.
}
\end{ex}

\begin{ex} %[2D1H03]
Tính hiệu biết số bị trừ là $76$, số trừ là $32$.
\shortans{44}
\loigiai{
\textbf{Đáp số:} $44$\\
$76 - 32 = 44$.
}
\end{ex}

\begin{ex} %[2D1V03]
Tính tổng của số tròn chục bé nhất có hai chữ số và số lớn nhất có hai chữ số.
\shortans{109}
\loigiai{
\textbf{Đáp số:} $109$\\
$10 + 99 = 109$.
}
\end{ex}

\begin{ex} %[2D1V03]
Tính hiệu của số lớn nhất có hai chữ số và số bé nhất có hai chữ số.
\shortans{89}
\loigiai{
\textbf{Đáp số:} $89$\\
$99 - 10 = 89$.
}
\end{ex}

\begin{ex} %[2D1V03]
Tìm tổng của ba số hạng: $10, 20, 30$.
\shortans{60}
\loigiai{
\textbf{Đáp số:} $60$\\
$10 + 20 + 30 = 60$.
}
\end{ex}

\begin{ex} %[2D1C03]
Trong một phép trừ, số bị trừ là $50$, hiệu là $20$. Hỏi số trừ là bao nhiêu?
\shortans{30}
\loigiai{
\textbf{Đáp số:} $30$\\
Số trừ: $50 - 20 = 30$.
}
\end{ex}
